problem:  
Let \( (M^n, g_0) \) be a compact, connected, smooth Riemannian manifold with nonempty boundary \( \partial M \), \( n \ge 3 \).  
Assume that:  
- the interior curvature operator of \(g_0\) is **positive**, and  
- the boundary \( \partial M \) is **strictly convex**, i.e. all eigenvalues of its shape operator \(S_0\) (with respect to the inward normal) are positive.

Consider the **Ricci flow with Dirichlet boundary condition**  
\[
\frac{\partial g}{\partial t} = -2\,\mathrm{Ric}(g), \qquad g|_{\partial M} = g_0|_{\partial M}.
\]
Such a flow is known to exist for a short time under suitable compatibility conditions.

**Question (main):**  
Does the Dirichlet Ricci flow preserve strict boundary convexity, that is, if \(S_0>0\), does one necessarily have \(S_t>0\) for as long as the flow exists?  
In other words, is boundary convexity a parabolic invariant under the Ricci flow when the boundary metric is fixed?

**Further motivation (secondary question):**  
If convexity is preserved for all time and the normalized flow (rescaled to keep total volume constant)
\[
\frac{\partial \tilde g}{\partial t}
   = -2\,\mathrm{Ric}(\tilde g)
     + \tfrac{2}{n}r(\tilde g)\,\tilde g, 
   \qquad \tilde g|_{\partial M}=g_0|_{\partial M},
\]
exists globally, can its limit (if any) be an **Einstein filling** of the prescribed boundary metric with **constant mean curvature boundary**?

Why is it a "good" problem:  
The principal question—preservation of strict boundary convexity under Ricci flow with Dirichlet data—is precise, well-posed, and genuinely open. Analytically, while the Dirichlet boundary condition provides short-time existence, no general mechanism ensures that positivity of the boundary shape operator persists: Hamilton’s tensor maximum principle is an interior tool and does not control evolution at the boundary. Demonstrating convexity preservation would require developing new **boundary tensor maximum principles** coupling the boundary shape operator’s evolution to interior curvature evolution—a mathematically deep challenge bridging PDE analysis and Riemannian geometry.

Geometrically, the problem asks whether Ricci curvature diffusion respects a natural extrinsic condition on the boundary. Success would establish a new invariant geometric class under Ricci flow (“strictly convex Dirichlet boundaries”), while failure would indicate fundamental obstructions to controlling extrinsic geometry through intrinsic curvature evolution.

The secondary question places this analytic issue in the broader context of **Einstein fillability** of a given boundary metric—a geometric motivation rather than a necessary part of the well-posed core.  Overall, the problem is simple to state yet profoundly difficult, connecting geometric flows, boundary curvature analysis, and the Einstein/CMC boundary problem, thereby exemplifying a “good” research-level mathematical problem.